% ============================================================
%  ENTREGABLE HITO 2 — Consolidación, Validación y Despliegue
%  Ingeniería de Software II
% ============================================================
\documentclass[11pt, a4paper]{article}

% ── Codificación y lenguaje ──────────────────────────────────
\usepackage[utf8]{inputenc}
\usepackage[T1]{fontenc}
\usepackage[spanish, es-tabla]{babel}

% ── Tipografía ───────────────────────────────────────────────
\usepackage{lmodern}
\usepackage{microtype}

% ── Geometría ────────────────────────────────────────────────
\usepackage[top=2.5cm, bottom=2.5cm, left=3cm, right=3cm]{geometry}

% ── Color ────────────────────────────────────────────────────
\usepackage[table]{xcolor}
\definecolor{primary}{RGB}{31, 73, 125}
\definecolor{accent}{RGB}{192, 0, 0}
\definecolor{success}{RGB}{0, 112, 60}
\definecolor{rowgray}{RGB}{240, 243, 248}
\definecolor{rowdark}{RGB}{31, 73, 125}

% ── Tablas ───────────────────────────────────────────────────
\usepackage{tabularx}
\usepackage{booktabs}
\usepackage{array}
\usepackage{ltablex}

\usepackage{longtable}
\usepackage{textcomp}
\usepackage{amssymb}
\renewcommand{\arraystretch}{1.35}

% ── Listas ───────────────────────────────────────────────────
\usepackage{enumitem}
\setlist[itemize]{itemsep=3pt, topsep=4pt, leftmargin=1.6em}
\setlist[enumerate]{itemsep=3pt, topsep=4pt, leftmargin=1.8em}

% ── Secciones ────────────────────────────────────────────────
\usepackage{titlesec}
\titleformat{\section}
  {\large\bfseries\color{primary}}
  {\thesection.}{0.6em}{}
  [\vspace{-0.4em}\textcolor{primary}{\rule{\linewidth}{0.7pt}}\vspace{0.2em}]

\titleformat{\subsection}
  {\normalsize\bfseries\color{primary}}
  {\thesubsection.}{0.5em}{}

\titleformat{\subsubsection}
  {\normalsize\itshape\color{primary}}
  {}{0pt}{}

% ── Encabezado y pie ─────────────────────────────────────────
\usepackage{fancyhdr}
\pagestyle{fancy}
\fancyhf{}
\renewcommand{\headrulewidth}{0.4pt}
\fancyhead[L]{\small\color{gray} Hito 2 — Ingeniería de Software II}
\fancyhead[R]{\small\color{gray} \textit{Nombre del equipo}}
\fancyfoot[C]{\small\thepage}

% ── Hipervínculos ────────────────────────────────────────────
\usepackage[hidelinks]{hyperref}

% ── Espacio entre párrafos ───────────────────────────────────
\usepackage{parskip}
\setlength{\parskip}{6pt}

% ── Cajas de información ─────────────────────────────────────
\usepackage{tcolorbox}
\tcbuselibrary{skins, breakable}

\newtcolorbox{infobox}[1][]{
  colback=rowgray, colframe=primary,
  fonttitle=\bfseries\small\color{white},
  coltitle=white, attach boxed title to top left={yshift=-2mm, xshift=6mm},
  boxed title style={colback=primary, rounded corners},
  rounded corners, left=6pt, right=6pt, top=10pt, bottom=6pt,
  #1
}

\newtcolorbox{warnbox}[1][]{
  colback=red!5, colframe=accent,
  fonttitle=\bfseries\small,
  rounded corners, left=6pt, right=6pt, top=6pt, bottom=6pt,
  #1
}

\newtcolorbox{notebox}[1][]{
  colback=success!5, colframe=success,
  rounded corners, left=6pt, right=6pt, top=6pt, bottom=6pt,
  #1
}

% ── Gráficos ─────────────────────────────────────────────────
\usepackage{graphicx}
\usepackage{float}

% ────────────────────────────────────────────────────────────
%  DATOS DEL EQUIPO (editar aquí)
% ────────────────────────────────────────────────────────────
\newcommand{\nombreequipo}{Cleanpool}
\newcommand{\nombreproyecto}{CleanPool App}
\newcommand{\integrantes}{%
  Vicente Araya (PM) \quad|\quad
  Diego Carmona (SM) \quad|\quad
  Daniel Cortez \quad|\quad
  Constanza Malebran
}
\newcommand{\seccion}{300}
\newcommand{\fechaentrega}{29 de abril de 2026}

% ============================================================
\begin{document}
% ============================================================

% ────────────────── PORTADA ─────────────────────────────────
\thispagestyle{empty}
\begin{center}
  \vspace*{1.5cm}

  {\small\color{gray} Universidad Andrés Bello \quad·\quad Ingeniería de Software II}

  \vspace{1cm}
  \textcolor{primary}{\rule{\linewidth}{1.5pt}}
  \vspace{0.4cm}

  {\Huge\bfseries\color{primary} Hito 2}\\[0.5em]
  {\LARGE Consolidación, Validación y\\Despliegue en Producción}

  \vspace{0.4cm}
  \textcolor{primary}{\rule{\linewidth}{1.5pt}}
  \vspace{1.2cm}

  {\Large\bfseries \nombreproyecto}

  \vspace{1cm}

  \begin{tabular}{rl}
    \textbf{Equipo:}    & \nombreequipo \\[4pt]
    \textbf{Sección:}   & \seccion \\[4pt]
    \textbf{Entrega:}   & \fechaentrega \\[4pt]
    \textbf{Integrantes:} & \parbox[t]{9cm}{\integrantes}
  \end{tabular}

  \vfill

  {\small\color{gray} Documento de entrega oficial — semestre activo}
\end{center}

\newpage

% ────────────────── ÍNDICE ──────────────────────────────────
\tableofcontents
\newpage

% ============================================================
\section{Exposición del Problema y Propuesta de Valor}
% ============================================================

\subsection{Definición del Problema}

\begin{infobox}[title=Instrucción]
Describir el problema central que el sistema busca resolver. Explicar por qué existe, a quién afecta y cuál es su magnitud o frecuencia. Ser concreto y evitar definiciones genéricas.
\end{infobox}

\vspace{0.3em}
La gestión de la calidad del agua en piscinas residenciales y de pequeña escala suele depender de mediciones manuales esporádicas, lo que dificulta detectar desviaciones de pH, cloro o temperatura a tiempo. Esto afecta a los dueños porque incrementa el riesgo sanitario, el sobreuso de químicos y los costos de mantención. El sistema aborda este problema mediante monitoreo continuo con sensores IoT y recomendaciones accionables desde una aplicación cliente.

\subsection{Contexto y Usuarios Objetivo}

\begin{itemize}
  \item \textbf{Usuarios principales:} Propietarios particulares de piscinas que desean mantener el agua de sus piscinas en perfectas condiciones, pero que no poseen conocimientos técnicos acerca del mantenimiento en piscinas.
  \item \textbf{Entorno de uso:} Aplicación Flutter multiplataforma conectada a una API FastAPI y MongoDB Atlas, con ingesta de sensores desde dispositivos ESP mediante endpoints protegidos por API key.
  \item \textbf{Necesidad identificada:} Contar con visibilidad continua del estado del agua y soporte de decisiones para tratamiento manual, reduciendo dependencia de chequeos reactivos.
\end{itemize}

\subsection{Propuesta de Valor y Evidencia}

\begin{tabularx}{\textwidth}{@{} l X @{}}
  \toprule
  \textbf{Elemento}        & \textbf{Descripción} \\
  \midrule
  Caso de uso 1            & Usuario autenticado consulta estado de su piscina (pH, cloro, temperatura) a partir de lecturas recientes ingresadas manualmente. \\
  Caso de uso 2            & Usuario registra tratamiento manual y revisa historial de mantenciones para trazabilidad operativa. \\
  Caso de uso 3            & Usuario obtiene temperatura de su piscina en tiempo real a través de sensor de temperatura. \\
  Impacto esperado         & El usuario puede mantener en óptimas condiciones el agua de su piscina, además de mejorar la consistencia en decisiones de mantención. \\
  Validación preliminar    & El usuario obtiene temperatura de su piscina y además consigue conocimientos necesarios para la mantención de su piscina. \\
  \bottomrule
\end{tabularx}

% ============================================================
\section{Requerimientos Funcionales y No Funcionales}
% ============================================================

\subsection{Requerimientos Funcionales}

\begin{longtable}{@{} p{1.4cm} p{3.5cm} X p{2cm} @{}}
  \toprule
  \textbf{ID} & \textbf{Nombre} & \textbf{Descripción} & \textbf{Prioridad} \\
  \midrule
  \endfirsthead
  \toprule
  \textbf{ID} & \textbf{Nombre} & \textbf{Descripción} & \textbf{Prioridad} \\
  \midrule
  \endhead
  \multicolumn{4}{@{}l}{\textbf{Acceso y seguridad de usuarios}} \\
  RF-01 & Autenticación de usuarios & El sistema debe permitir registro, inicio de sesión y consulta de perfil autenticado mediante JWT. & Alta \\
  RF-15 & Ingreso seguro propietario & El sistema debe implementar autenticación segura para el ingreso del propietario de la piscina. & Alta \\
  \midrule
  \multicolumn{4}{@{}l}{\textbf{Gestión de piscinas y dispositivos IoT}} \\
  RF-02 & Gestión de piscinas & El sistema debe permitir crear, listar, actualizar y eliminar piscinas asociadas al usuario autenticado. & Alta \\
  RF-05 & Ingesta y vínculo de dispositivos & El sistema debe recibir telemetría IoT protegida por API key y permitir vincular/desvincular dispositivos a piscinas. & Media \\
  RF-13 & Estado al vincular IoT & El sistema debe mostrar el estado del agua cada vez que se vincule el dispositivo (IoT) a la app. & Media \\
  \midrule
  \multicolumn{4}{@{}l}{\textbf{Medición, evaluación y visualización del estado del agua}} \\
  RF-06 & Medición de temperatura & El sistema debe medir temperatura del agua. & Alta \\
  RF-07 & Comparación con rangos óptimos & El sistema debe comparar automáticamente los valores medidos con rangos óptimos predefinidos para piscinas residenciales. & Alta \\
  RF-08 & Indicadores fuera de rango & El sistema debe generar indicadores cuando uno o más parámetros estén fuera de rango seguro, indicando el parámetro afectado con su valor. & Alta \\
  RF-03 & Monitoreo de estado & El sistema debe evaluar lecturas de sensores y entregar estado por parámetro y estado global de aptitud. & Alta \\
  RF-12 & Aptitud para baño & El sistema debe mostrar en el dashboard si el agua está apta o no apta para el baño, basándose en la evaluación conjunta de los parámetros medidos. & Alta \\
  RF-11 & Dashboard móvil & El sistema debe permitir al usuario consultar el estado actual de la piscina desde un dashboard dentro de una app móvil, accesible desde cualquier smartphone. & Alta \\
  RF-10 & Historial de mediciones & El sistema debe almacenar el historial de mediciones con identificador de sensor y valor registrado. & Alta \\
  \midrule
  \multicolumn{4}{@{}l}{\textbf{Recomendación y acciones de mantención}} \\
  RF-09 & Recomendación química & El sistema debe recomendar qué producto químico aplicar y en qué cantidad aproximada según volumen de piscina y valores actuales. & Alta \\
  RF-14 & Acción recomendada & El sistema debe mostrar la acción recomendada en base a los parámetros del agua. & Alta \\
  RF-04 & Mantenciones y tratamientos & El sistema debe calcular dosis de tratamiento manual y persistir historial de mantenciones. & Media \\
  \bottomrule
\end{longtable}

\subsection{Requerimientos No Funcionales}

\begin{longtable}{@{} p{1.4cm} p{2.8cm} X p{2.2cm} @{}}
  \toprule
  \textbf{ID} & \textbf{Categoría} & \textbf{Descripción} & \textbf{Métrica} \\
  \midrule
  \endfirsthead
  \toprule
  \textbf{ID} & \textbf{Categoría} & \textbf{Descripción} & \textbf{Métrica} \\
  \midrule
  \endhead
  \multicolumn{4}{@{}l}{\textbf{Seguridad}} \\
  RNF-03 & Seguridad & Las credenciales deben almacenarse hasheadas y el acceso a endpoints protegidos debe validarse con token Bearer JWT; la ingesta IoT exige API key. & JWT + bcrypt + X-API-KEY \\
  RNF-10 & Seguridad & El acceso a los datos debe requerir autenticación con token JWT. & JWT obligatorio \\
  \midrule
  \multicolumn{4}{@{}l}{\textbf{Disponibilidad e infraestructura}} \\
  RNF-02 & Disponibilidad & El backend debe ser desplegable en entorno productivo y ejecutable localmente por Docker o Uvicorn.\\
  RNF-07 & Disponibilidad & El sistema debe estar continuamente disponible un 90\% del tiempo a nivel semanal garantizando su estabilidad. & 90\% semanal \\
  RNF-15 & Disponibilidad en nube & El sistema debe estar desplegado en infraestructura cloud que garantice una disponibilidad mínima del 97\% mensual, excluyendo ventanas de mantenimiento programadas. & 97\% mensual \\
  \midrule
  \multicolumn{4}{@{}l}{\textbf{Rendimiento y latencia}} \\
  RNF-01 & Rendimiento & La API debe responder de forma interactiva para operaciones de consulta y registro sobre lecturas, piscinas y auth.\\
  RNF-09 & Latencia de actualización & La visualización de datos en el dashboard debe actualizarse en al menos 5 segundos desde la captura del sensor. & $\leq 5$ segundos \\
  RNF-13 & Rendimiento funcional & El sistema debe mostrar recomendaciones de productos químicos en menos de 5 segundos una vez que se hace clic en solicitar la recomendación. & $\leq 5$ segundos \\
  \midrule
  \multicolumn{4}{@{}l}{\textbf{Calidad de datos y precisión}} \\
  RNF-08 & Precisión & Las mediciones de sensores tienen un margen de error menor a $\pm$1 en temperatura. & Error $< \pm1\,^{\circ}\mathrm{C}$ \\
  RNF-12 & Validación de datos & El sistema debe validar cada lectura antes del almacenamiento, aceptando solo pH entre 0.0 y 14.0, temperatura entre 0$^\circ$C y 60$^\circ$C y cloro entre 0.0 y 10.0 ppm; las lecturas fuera de rango se descartan.\\
  \midrule
  \multicolumn{4}{@{}l}{\textbf{Usabilidad y compatibilidad}} \\
  RNF-04 & Usabilidad & La interfaz Flutter debe presentar navegación por pantallas de autenticación, dashboard, perfil y gestión de piscina con flujo comprensible para usuario final. \\
  RNF-06 & Usabilidad & La interfaz debe mostrar los niveles de pH, cloro y temperatura al hacer clic en dashboard. & Visualización en dashboard \\
  RNF-14 & Compatibilidad móvil & La app debe ser completamente funcional en dispositivos iOS versión 14.0 o superior y Android 10.0 o superior. & iOS 14+ y Android 10+ \\
  \midrule
  \multicolumn{4}{@{}l}{\textbf{Mantenibilidad y evolución}} \\
  RNF-11 & Mantenibilidad & El sistema debe diseñarse de forma modular con capas separadas en IoT, lógica de negocio, datos y presentación para facilitar correcciones y mejoras sin afectar otros módulos. & Arquitectura en capas \\
  RNF-05 & Escalabilidad & La arquitectura desacoplada frontend/backend y uso de MongoDB Atlas permite crecer en clientes y lecturas sin acoplar UI con almacenamiento. \\
  \bottomrule
\end{longtable}

% ============================================================
\section{Diseño y Justificación Arquitectónica}
% ============================================================

\subsection{Arquitectura del Sistema}

% \begin{figure}[H]
%   \centering
%   \includegraphics[width=0.85\textwidth]{arquitectura.png}
%   \caption{Diagrama de arquitectura del sistema — Hito 2}
% \end{figure}

La arquitectura se compone de: (1) cliente Flutter (móvil) como capa de presentación; (2) API REST FastAPI como capa de aplicación; (3) MongoDB Atlas como persistencia; y (4) integración de ingesta IoT mediante endpoints versionados con autenticación por API key para dispositivos. El backend expone además autenticación JWT para usuarios finales y servicios de negocio para evaluación de estado y mantenciones.

\subsection{Cambios Respecto al Hito Anterior}

\begin{tabularx}{\textwidth}{@{} l X X @{}}
  \toprule
  \textbf{Componente} & \textbf{Estado anterior} & \textbf{Estado actual} \\
  \midrule
  Backend API & Endpoints base de autenticación/lecturas y estructura inicial. & Se consolidan routers para ingesta IoT, vinculación de dispositivos, piscinas de usuario y mantenciones. \\
  Cliente Flutter & Flujo inicial de interfaz y consumo básico de API. & Se estructura por features con servicios compartidos para auth, pool y mantenciones, incluyendo persistencia local de sesión. \\
  Placa arduino & Uso de arduino ESP32 & Uso de arduino ESP8266 \\
  Access point desde esp82666 para conexión con la app & esp8266 se conecta a red Wi-fi y asi manda datos con protocolo HTTP al backend de la app \\
  \bottomrule
\end{tabularx}

\subsection{Decisiones de Diseño y Trade-offs}

\begin{longtable}{@{} p{3.5cm} X X @{}}
  \toprule
  \textbf{Decisión} & \textbf{Justificación} & \textbf{Trade-off asumido} \\
  \midrule
  \endfirsthead
  \textbf{Decisión} & \textbf{Justificación} & \textbf{Trade-off asumido} \\
  \midrule
  \endhead
  REST (FastAPI) vs GraphQL & Menor complejidad operativa y rapidez de implementación para endpoints orientados a recursos del dominio. & Simplicidad y estandarización HTTP, a cambio de menor flexibilidad en agregación de consultas complejas. \\
  Monolito modular vs microservicios & Tamaño del equipo y del dominio favorecen una base única con routers/servicios separados por responsabilidad. & Menor costo de despliegue y coordinación, a cambio de mayor riesgo de crecimiento acoplado si no se controla deuda técnica. \\
  \bottomrule
\end{longtable}

\subsection{Problemas Detectados y Soluciones}

\begin{tabularx}{\textwidth}{@{} p{4cm} X @{}}
  \toprule
  \textbf{Problema detectado} & \textbf{Solución implementada} \\
  \midrule
  No se pudo hacer funcionar el sensor de ph ya que era imposible de calibrar. & Se mantiene sensor de temperatura funcionando\\
  Inconsistencia entre rutas históricas y rutas actuales en pruebas heredadas. & Se mantiene cobertura principal con tests específicos por módulo (pools, mantenciones, status) y se documenta el desfase para normalizar endpoints. \\
  El arduino no lograba enviar los datos del sensor de temperatura a la aplicación & Se modificó el endpoint necesario apuntando al backend desplegado para mostrar temperatura en aplicación\\
  El modelado 3D del diseño al principio generó problemas de impresión que hicieron que no se pudiera imprimir por ancho de paredes & Se ajustaron las medidas del modelo del flotador en tinkercad para ser impreso en 3D\\
  \bottomrule
\end{tabularx}

% ============================================================
\section{Estrategia de Pruebas}
% ============================================================

\subsection{Enfoque General}

Se adopta un enfoque mixto con pruebas unitarias y de integración en backend usando Pytest, FastAPI TestClient y mocks de MongoDB para desacoplar dependencias externas. En frontend existe prueba base de widget para validar arranque. La priorización se centra en autenticación, estado de piscina, mantenciones y configuración de pools.

\subsection{Plan de Pruebas}

\begin{longtable}{@{} p{1.3cm} p{2.5cm} p{2.5cm} X p{2cm} @{}}
  \toprule
  \textbf{ID} & \textbf{Tipo} & \textbf{Módulo/Función} & \textbf{Descripción} & \textbf{Resultado} \\
  \midrule
  \endfirsthead
  \toprule
  \textbf{ID} & \textbf{Tipo} & \textbf{Módulo/Función} & \textbf{Descripción} & \textbf{Resultado} \\
  \midrule
  \endhead
  TP-01 & Unitaria       & status\_service & Valida cálculo de estado global de aptitud desde estados de parámetros. & \textit{[Implementada en repositorio]} \\
  TP-02 & Unitaria       & pools/settings & Verifica comportamiento de defaults y validaciones de configuración de piscina en modelos y endpoints de pools. & \textit{[Implementada en repositorio]} \\
  TP-03 & Integración    & auth + mantenciones & Comprueba flujo autenticado para crear/listar mantenciones y restricciones sin token. & \textit{[Implementada en repositorio]} \\
  TP-04 & Integración    & test\_api legado & Incluye casos sobre /readings que no coinciden completamente con routers activos en main. & \textit{[Requiere alineación]} \\
  TP-05 & Sistema  & frontend widget test  & Valida construcción básica de la app Flutter en entorno de pruebas. & \textit{[Implementada en repositorio]} \\
  TP-06 & Unitaria       & config\_pool/evaluar\_sensor & Evalúa rangos y clasificación de estado (OPTIMO, ADVERTENCIA, CRITICO), incluyendo error por sensor no reconocido. & \textit{[Implementada en repositorio]} \\
  TP-07 & Unitaria       & services/calculator & Valida aptitud global simple, estado individual, validación de volumen y reglas combinadas de dosificación (orden pH $\rightarrow$ cloro). & \textit{[Implementada en repositorio]} \\
  TP-08 & Unitaria       & config/Settings & Verifica validación de SECRET\_KEY y aplicación de valores por defecto de configuración. & \textit{[Implementada en repositorio]} \\
  TP-09 & Unitaria       & routers/ingesta & Cubre verify\_api\_key, persistencia de is\_critical y rechazo de payload legacy con pH fuera de rango físico. & \textit{[Implementada en repositorio]} \\
  TP-10 & Unitaria       & routers/device\_bindings & Cubre vínculo/desvínculo y estado de dispositivo (online/offline), además de ingesta por device\_id con y sin binding activo. & \textit{[Implementada en repositorio]} \\
  \bottomrule
\end{longtable}

\subsection{Cobertura y Herramientas}

\begin{itemize}
  \item \textbf{Herramientas utilizadas:} Pytest, FastAPI TestClient, unittest.mock (backend) y flutter\_test (frontend).
  \item \textbf{Cobertura de código alcanzada:} Se añadieron 26 pruebas unitarias nuevas para backend (config\_pool, services/calculator, config/settings, routers/ingesta y routers/device\_bindings), todas en estado \textit{pass} al ejecutarse de forma aislada.
  \item \textbf{Criterio de aceptación:} Los módulos críticos (auth, estado de piscina, mantenciones, pools) y módulos base de lógica/ingesta deben contar con pruebas automatizadas; las pruebas legacy desalineadas deben normalizarse con los endpoints vigentes.
\end{itemize}

% ============================================================
\section{Identificación y Gestión de Riesgos}
% ============================================================

\subsection{Matriz de Riesgos}

\begin{longtable}{@{} p{1.2cm} X p{1.8cm} p{1.6cm} p{2.8cm} p{2cm} @{}}
  \toprule
  \textbf{ID} & \textbf{Descripción del riesgo} & \textbf{Probabilidad} & \textbf{Impacto} & \textbf{Plan de mitigación} & \textbf{Estado} \\
  \midrule
  \endfirsthead
  \toprule
  \textbf{ID} & \textbf{Descripción del riesgo} & \textbf{Probabilidad} & \textbf{Impacto} & \textbf{Plan de mitigación} & \textbf{Estado} \\
  \midrule
  \endhead
  R-01 & Dependencia de servicios externos (MongoDB Atlas / nube) para operación continua. & Media   & Alto   & Mantener variables de entorno claras, validación de conexión y plan de contingencia para entorno alternativo/local. & \textcolor{accent}{Activo} \\
  R-02 & Desfase entre rutas legacy y rutas vigentes en pruebas/documentación.  & Media  & Medio   & Homologar contratos API, depurar endpoints en desuso y actualizar suite de pruebas.       & \textcolor{accent}{Activo} \\
  R-03 & CORS completamente abierto en backend en un eventual despliegue público.     & Media   & Medio  & Restringir orígenes permitidos y endurecer configuración por entorno.      & \textcolor{accent}{Activo} \\
  R-04 & Deuda técnica por duplicidad de modelos/colecciones (piscinas/pools, mantenciones/mantenimientos).          & Media  & Medio  & Planificar refactor de unificación de dominio y normalización de repositorios. & Activo \\
  \bottomrule
\end{longtable}

\subsection{Riesgos Materializados y Respuesta}

Se observa materialización parcial del riesgo de inconsistencia técnica entre componentes legacy y actuales (por ejemplo, rutas probadas históricamente versus routers incluidos en la aplicación). El impacto principal es aumento de costo de mantenimiento y ambigüedad documental. Como respuesta, se prioriza normalización de contratos API, consolidación de nomenclatura de colecciones y actualización de pruebas de integración.

% ============================================================
\section{Incremento Funcional y Despliegue}
% ============================================================

\subsection{Funcionalidades Implementadas en este Hito}

\begin{longtable}{@{} p{1.4cm} p{4cm} X p{2cm} @{}}
  \toprule
  \textbf{ID} & \textbf{Historia de usuario} & \textbf{Descripción} & \textbf{Estado} \\
  \midrule
  \endfirsthead
  \toprule
  \textbf{ID} & \textbf{Historia de usuario} & \textbf{Descripción} & \textbf{Estado} \\
  \midrule
  \endhead
  HU-01 & Como usuario, quiero registrarme e iniciar sesión para acceder de forma segura a mis datos de piscina. & Implementación en endpoints /auth/register, /auth/login, /auth/me y persistencia local de sesión en Flutter. & \textcolor{success}{\textbf{Completa}} \\
  HU-02 & Como usuario, quiero visualizar el estado de mi piscina para decidir acciones de mantención. & Evaluación de lecturas y estado global, además de endpoints de consulta y pantallas de dashboard/perfil. & \textcolor{success}{\textbf{Completa}} \\
  HU-03 & Como usuario, quiero que el dispositivo quede vinculado de forma trazable a mi piscina para automatizar lecturas. & Existen endpoints de bind/unbind/status, pero persisten elementos legacy y normalización pendiente entre modelos de dominio. & \textcolor{accent}{\textbf{Parcial}} \\
  \bottomrule
\end{longtable}

\subsection{Entorno de Despliegue}

\begin{tabularx}{\textwidth}{@{} l X @{}}
  \toprule
  \textbf{Parámetro}          & \textbf{Detalle} \\
  \midrule
  Plataforma                  & Azure Web App (backend) y despliegue frontend web por configuración Docker/Nginx y Vercel SPA config en repositorio. \\
  URL de acceso               & \textit{[No visible explícitamente en el repositorio]} \\
  Credenciales de prueba      & \textit{[No deben publicarse; no hay credenciales de demo explícitas en repositorio]} \\
  Estado al momento de entrega & Backend con pipeline de despliegue configurado y ejecución local documentada para backend/frontend. \\
  Notas adicionales           & Endpoints de usuario requieren Bearer JWT; ingesta de sensores requiere header X-API-KEY. \\
  \bottomrule
\end{tabularx}

\subsection{Evidencia de Funcionamiento}

\begin{infobox}[title=Instrucción]
Incluir capturas de pantalla o descripciones de las funcionalidades clave en el entorno de producción. Si se incluyen imágenes, descomentar el bloque \texttt{figure} y reemplazar el nombre del archivo.
\end{infobox}

% \begin{figure}[H]
%   \centering
%   \includegraphics[width=0.8\textwidth]{captura1.png}
%   \caption{Descripción de la funcionalidad mostrada.}
% \end{figure}

Se puede evidenciar el flujo completo: registro/login, visualización de estado de piscina en dashboard, ejecución de tratamiento manual y revisión de historial de mantenciones, además del registro de datos de sensores en backend vía endpoints de ingesta.

% ============================================================
\section{Aplicación de la Metodología Scrum}
% ============================================================

\subsection{Gestión del Sprint (Product Manager)}

\begin{tabularx}{\textwidth}{@{} l X @{}}
  \toprule
  \textbf{Parámetro}            & \textbf{Detalle} \\
  \midrule
  Número de sprint              & \textit{[Sprint \#X]} \\
  Objetivo del sprint           & \textit{[Descripción del objetivo comprometido.]} \\
  Grado de cumplimiento         & \textit{[ej. 80\% — 4 de 5 historias completadas.]} \\
  Ajustes realizados            & \textit{[Cambios en el alcance o prioridad y su justificación.]} \\
  Objetivo del siguiente sprint & \textit{[Descripción del sprint planificado.]} \\
  \bottomrule
\end{tabularx}

\vspace{0.5em}

\subsubsection{Estado del Product Backlog}

\begin{longtable}{@{} p{1.3cm} X p{2.5cm} p{2cm} @{}}
  \toprule
  \textbf{ID} & \textbf{Historia de usuario} & \textbf{Sprint asignado} & \textbf{Estado} \\
  \midrule
  \endfirsthead
  \toprule
  \textbf{ID} & \textbf{Historia de usuario} & \textbf{Sprint asignado} & \textbf{Estado} \\
  \midrule
  \endhead
  HU-01 & \textit{[descripción]} & Sprint 1 & \textcolor{success}{\textbf{Completa}} \\
  HU-02 & \textit{[descripción]} & Sprint 2 & \textcolor{success}{\textbf{Completa}} \\
  HU-03 & \textit{[descripción]} & Sprint 2 & \textcolor{accent}{\textbf{Parcial}} \\
  HU-04 & \textit{[descripción]} & Sprint 3 & Pendiente \\
  HU-05 & \textit{[descripción]} & Sprint 3 & Pendiente \\
  \bottomrule
\end{longtable}

\subsection{Ejecución Scrum (Scrum Master)}

\subsubsection{Daily Meetings}

\textit{[Describir la dinámica de los Daily Meetings: frecuencia, modalidad (presencial/remota), duración promedio y disciplina del equipo.]}

\subsubsection{Sprint Retrospective — Principales Mejoras}

\begin{tabularx}{\textwidth}{@{} p{3.5cm} X @{}}
  \toprule
  \textbf{Aspecto} & \textbf{Acción tomada} \\
  \midrule
  \textit{[Problema identificado]} & \textit{[Mejora implementada]} \\
  \textit{[Problema identificado]} & \textit{[Mejora implementada]} \\
  \bottomrule
\end{tabularx}

\subsubsection{Impedimentos y Adaptaciones del Equipo}

\textit{[Describir los principales impedimentos encontrados durante el sprint y cómo el equipo los resolvió o mitigó.]}

\subsection{Registro de Asistencia a Daily Meetings}

\begin{warnbox}
  \textbf{Nota:} Este componente tiene alto peso evaluativo. El registro debe ser completo, individual y reflejar la participación real del equipo.
\end{warnbox}

\vspace{0.4em}

\begin{tabularx}{\textwidth}{|l|*{6}{>{\centering\arraybackslash}p{1.1cm}|}|c|c|}
\hline
\rowcolor{rowdark}
\textcolor{white}{\textbf{Integrante}} &
\textcolor{white}{\textbf{D1}} &
\textcolor{white}{\textbf{D2}} &
\textcolor{white}{\textbf{D3}} &
\textcolor{white}{\textbf{D4}} &
\textcolor{white}{\textbf{D5}} &
\textcolor{white}{\textbf{D6}} &
\textcolor{white}{\textbf{Asistidas}} &
\textcolor{white}{\textbf{\%}} \\
\hline
\textit{Integrante 1 (PM)} & \textcolor{success}{\checkmark} & \textcolor{success}{\checkmark} & \textcolor{success}{\checkmark} & \textcolor{accent}{$\times$} & \textcolor{success}{\checkmark} & \textcolor{success}{\checkmark} & 5/6 & 83\% \\
\hline
\textit{Integrante 2 (SM)} & \textcolor{success}{\checkmark} & \textcolor{success}{\checkmark} & \textcolor{success}{\checkmark} & \textcolor{success}{\checkmark} & \textcolor{success}{\checkmark} & \textcolor{success}{\checkmark} & 6/6 & 100\% \\
\hline
\textit{Integrante 3}      & \textcolor{success}{\checkmark} & \textcolor{accent}{$\times$} & \textcolor{success}{\checkmark} & \textcolor{success}{\checkmark} & \textcolor{accent}{$\times$} & \textcolor{success}{\checkmark} & 4/6 & 67\% \\
\hline
\textit{Integrante 4}      & \textcolor{success}{\checkmark} & \textcolor{success}{\checkmark} & \textcolor{success}{\checkmark} & \textcolor{success}{\checkmark} & \textcolor{success}{\checkmark} & \textcolor{accent}{$\times$} & 5/6 & 83\% \\
\hline
\end{tabularx}

\vspace{0.3em}

\textit{Justificación de inasistencias:}
\begin{itemize}
  \item \textbf{Integrante 1, D4:} \textit{[Motivo justificado.]}
  \item \textbf{Integrante 3, D2 y D5:} \textit{[Motivo justificado.]}
\end{itemize}

% ============================================================
\section{Presentación Final Orientada a Stakeholders}
% ============================================================

\subsection{Narrativa del Proyecto}

CleanPool App aborda el problema de control reactivo de calidad de agua en piscinas, entregando monitoreo digital y apoyo a decisiones de mantención para usuarios responsables de operación. En este hito se consolidó una arquitectura operativa con frontend Flutter, backend FastAPI y persistencia en MongoDB Atlas, incorporando autenticación, gestión de piscinas, ingesta IoT, evaluación de estado y registro de mantenciones. El valor entregado es una base funcional desplegable que centraliza información crítica y mejora la trazabilidad de acciones técnicas.

\subsection{Avances Clave del Hito}

\begin{itemize}
  \item Se implementó flujo de autenticación completo con sesión persistente en cliente Flutter.
  \item Se habilitó monitoreo del estado del agua con evaluación por parámetros y estado global.
  \item Se incorporó gestión de mantenciones y tratamiento manual con historial consultable.
\end{itemize}

\subsection{Demostración en Producción}

\begin{notebox}
  El sistema se encuentra desplegado y disponible para evaluación en: \textbf{\textit{[URL no especificada en repositorio]}}. \\
  Credenciales de acceso para la demo: \textbf{Usuario:} \textit{[No disponible en repositorio]} \quad \textbf{Contraseña:} \textit{[No disponible en repositorio]}
\end{notebox}

\vspace{0.3em}
Flujo sugerido de demostración: (1) registro/inicio de sesión, (2) creación/selección de piscina, (3) revisión de estado actual y métricas, (4) ejecución de tratamiento manual, (5) validación de historial de mantenciones y (6) evidencia de endpoints de ingesta/vinculación de dispositivo.

\subsection{Próximos Pasos}

\begin{enumerate}
  \item Normalizar y unificar modelos de dominio duplicados entre rutas legacy y rutas vigentes.
  \item Alinear pruebas, documentación y contratos de API para eliminar rutas obsoletas y asegurar trazabilidad funcional.
  \item Endurecer configuración de seguridad y observabilidad (CORS por entorno, logging estructurado, reporte de cobertura).
\end{enumerate}

% ============================================================
\appendix
\renewcommand{\thesection}{Anexo \Alph{section}}
% ============================================================

\section{Capturas de Pantalla del Sistema}

\textit{[Pendiente adjuntar capturas reales del entorno desplegado.]}

\section{Actas de Daily Meetings}

\textit{[Pendiente adjuntar actas reales del equipo.]}

\section{Evidencia de Repositorio}

\begin{tabularx}{\textwidth}{@{} l X @{}}
  \toprule
  \textbf{Recurso} & \textbf{URL} \\
  \midrule
  Repositorio (GitHub / GitLab) & \textit{[URL exacta del repositorio no especificada en este archivo]} \\
  Tablero Scrum (Jira / Trello)  & \textit{[No evidenciado en repositorio]} \\
  Sistema desplegado             & \textit{[URL pública no evidenciada en repositorio]} \\
  \bottomrule
\end{tabularx}

% ============================================================
\end{document}
% ============================================================
